% =====================================================================
\chapter{Sucesiones y límites}
\label{tema:2}
% =====================================================================

El concepto central del análisis es el de límite. Este tema lo introduce primero para sucesiones, donde la definición es más sencilla, y después para funciones, con la definición $\varepsilon$--$\delta$. Se demuestran las reglas de cálculo de límites, el teorema del sándwich y el teorema de la convergencia monótona; se obtienen los límites notables que se usarán en el resto del módulo y se estudian las asíntotas. Se siguen los capítulos~5 y~22 de Spivak~\cite{spivak}, el capítulo~2 y la §5.1 de los apuntes de \emph{Analysis~I} de Cambridge~\cite{cambA1}, y las §2.2--2.5 y~4.6 de OpenStax~\cite{osc1}.

\section{Sucesiones convergentes}

\begin{definicion}[Sucesión y límite]\label{def:limite-sucesion}
Una \emph{sucesión} de números reales es una función $\N\to\R$, $n\mapsto a_n$; se escribe $(a_n)_{n\ge1}$ o $(a_n)$. Se dice que $(a_n)$ \emph{converge} a $L\in\R$, y se escribe $a_n\to L$ o $\lim_{n\to\infty}a_n=L$, si
\[
\forall\varepsilon>0\ \ \exists N\in\N\ \ \forall n\ge N:\quad \abs{a_n-L}<\varepsilon.
\]
Una sucesión que no converge a ningún número real es \emph{divergente}.
\end{definicion}

La definición se lee como un juego: un adversario fija una tolerancia $\varepsilon>0$, tan pequeña como quiera, y hay que responder con un índice $N$ a partir del cual todos los términos están a distancia menor que $\varepsilon$ de $L$. El índice $N$ puede depender de $\varepsilon$, y casi siempre depende.

\addtocounter{definicion}{1}%

\addtocounter{definicion}{1}%

En las demostraciones $\varepsilon$--$N$ conviene separar el trabajo previo (acotar $\abs{a_n-L}$ por una expresión sencilla que tienda a~0) de la demostración en sí, que se redacta en el orden de los cuantificadores: «sea $\varepsilon>0$; sea $N=\dots$; para $n\ge N$, \dots».

\begin{proposicion}[Unicidad y acotación]\label{prop:unicidad}
(i) El límite de una sucesión convergente es único. (ii) Toda sucesión convergente está acotada.
\end{proposicion}

\begin{teorema}[Álgebra de límites de sucesiones]\label{teo:algebra-sucesiones}
Sean $a_n\to L$ y $b_n\to M$. Entonces
(i) $a_n+b_n\to L+M$; (ii) $ca_n\to cL$ para todo $c\in\R$; (iii) $a_nb_n\to LM$; (iv) si $M\ne0$ y $b_n\ne0$ para todo $n$, $a_n/b_n\to L/M$.
\end{teorema}

\begin{teorema}[Teorema del sándwich]\label{teo:sandwich}
Si $a_n\le b_n\le c_n$ para todo $n\ge n_0$ y $a_n\to L$, $c_n\to L$, entonces $b_n\to L$.
\end{teorema}

\begin{definicion}[Divergencia a infinito]\label{def:div-infinito}
$a_n\to+\infty$ si para todo $K\in\R$ existe $N$ tal que $a_n>K$ para todo $n\ge N$. Análogamente, $a_n\to-\infty$.
\end{definicion}

\begin{proposicion}[Sucesión geométrica]\label{prop:geometrica}
Sea $r\in\R$. Si $\abs{r}<1$, $r^{n}\to0$; si $r>1$, $r^{n}\to+\infty$; si $r=1$, $r^{n}\to1$; si $r\le-1$, $(r^{n})$ diverge.
\end{proposicion}

\section{Sucesiones monótonas y el número \texorpdfstring{$e$}{e}}

\begin{teorema}[Convergencia monótona]\label{teo:monotona}
Toda sucesión creciente y acotada superiormente converge a $\sup\{a_n:n\in\N\}$. Toda sucesión decreciente y acotada inferiormente converge a su ínfimo.
\end{teorema}

Este teorema es la primera consecuencia importante de la completitud de $\R$: en $\Q$ es falso, porque la sucesión creciente y acotada $1;\ 1{,}4;\ 1{,}41;\ 1{,}414;\dots$ de aproximaciones decimales de $\sqrt2$ no tiene límite racional.

\begin{teorema}[El número $e$]\label{teo:e}
La sucesión $e_n=\bigl(1+\tfrac1n\bigr)^{n}$ es creciente y está acotada superiormente por~3. Su límite se denota $e$ y vale $e=2{,}71828\dots$
\end{teorema}

\section{Límite de una función}

\begin{definicion}[Límite $\varepsilon$--$\delta$]\label{def:limite-funcion}
Sea $f$ definida en un intervalo abierto que contiene a $a$, salvo quizá en el propio $a$. Se dice que $f$ tiene límite $L$ en $a$, y se escribe $\lim_{x\to a}f(x)=L$, si
\[
\forall\varepsilon>0\ \ \exists\delta>0\ \ \forall x:\quad 0<\abs{x-a}<\delta\ \Rightarrow\ \abs{f(x)-L}<\varepsilon.
\]
\end{definicion}

La condición $0<\abs{x-a}$ excluye el punto $a$: el límite describe el comportamiento de $f$ \emph{cerca} de $a$, no \emph{en} $a$, y no depende de $f(a)$ ni de que esté definido.

\addtocounter{definicion}{1}%

\addtocounter{definicion}{1}%

\begin{teorema}[Caracterización por sucesiones]\label{teo:heine}
$\lim_{x\to a}f(x)=L$ si y solo si, para toda sucesión $(x_n)$ con $x_n\ne a$ y $x_n\to a$, se tiene $f(x_n)\to L$.
\end{teorema}

El \cref{teo:heine} traslada a las funciones todo lo demostrado para sucesiones: el límite es único, y valen el álgebra de límites (\cref{teo:algebra-sucesiones}) y el teorema del sándwich para funciones. También sirve para probar que un límite \emph{no} existe: basta encontrar dos sucesiones $x_n\to a$ y $y_n\to a$ con $\lim f(x_n)\ne\lim f(y_n)$.

\addtocounter{definicion}{1}%

\begin{definicion}[Límites laterales, infinitos y en el infinito]\label{def:limites-variantes}
Sea $f$ una función real y $a\in\R$.
\begin{enumerate}
  \item $\lim_{x\to a^{+}}f(x)=L$ si para todo $\varepsilon>0$ existe $\delta>0$ con $\abs{f(x)-L}<\varepsilon$ siempre que $a<x<a+\delta$. Análogamente, $\lim_{x\to a^{-}}$ con $a-\delta<x<a$.
  \item $\lim_{x\to a}f(x)=+\infty$ si para todo $K$ existe $\delta>0$ con $f(x)>K$ siempre que $0<\abs{x-a}<\delta$.
  \item $\lim_{x\to+\infty}f(x)=L$ si para todo $\varepsilon>0$ existe $M$ con $\abs{f(x)-L}<\varepsilon$ siempre que $x>M$.
\end{enumerate}
\end{definicion}

El límite existe si y solo si existen los dos laterales y coinciden (\cref{eje:2-laterales}). Por ejemplo, para $f(x)=\abs{x}/x$ los laterales en~0 son $1$ y $-1$, y el límite no existe.

\section{Límites notables}

\begin{proposicion}[$\sin x/x$]\label{prop:senx}
Para $0<\abs{x}<\pi/2$ se cumple $\cos x<\dfrac{\sin x}{x}<1$. En consecuencia, $\displaystyle\lim_{x\to0}\frac{\sin x}{x}=1$ y $\displaystyle\lim_{x\to0}\frac{1-\cos x}{x^{2}}=\frac12$.
\end{proposicion}

Para la exponencial se adopta el punto de vista del curso 18.01 del MIT~\cite{mit1801} (sesiones~16 a~19): el número $e$ es la única base $b>1$ para la que la gráfica de $b^{x}$ tiene pendiente~1 en el origen. Se admite, por tanto, que la exponencial es continua y que
\begin{equation}\label{eq:lim-exp}
\lim_{h\to0}\frac{e^{h}-1}{h}=1,
\end{equation}
y que esta base coincide con el límite del \cref{teo:e}. En Spivak~\cite[cap.~18]{spivak} se construyen $\ln$ y $\exp$ a partir de la integral y todo lo anterior se demuestra.

\begin{proposicion}[Límites notables de la exponencial y el logaritmo]\label{prop:notables-exp}
\[
\text{(i)}\ \lim_{x\to0}\frac{\ln(1+x)}{x}=1;\qquad
\text{(ii)}\ \lim_{x\to0}(1+x)^{1/x}=e;\qquad
\text{(iii)}\ \lim_{n\to\infty}\Bigl(1+\frac{t}{n}\Bigr)^{n}=e^{t}\ \ (t\in\R).
\]
\end{proposicion}

\begin{table}[h]
\centering\small
\caption{Límites notables en $x\to0$. Se demuestran en este tema o en sus ejercicios.}\label{tab:notables}
\begin{tabular}{@{}lll@{}}
\toprule
$\dfrac{\sin x}{x}\to1$ & $\dfrac{1-\cos x}{x^{2}}\to\dfrac12$ & $\dfrac{\tan x}{x}\to1$ \\[2mm]
$\dfrac{e^{x}-1}{x}\to1$ & $\dfrac{\ln(1+x)}{x}\to1$ & $(1+x)^{1/x}\to e$ \\[2mm]
$\dfrac{a^{x}-1}{x}\to\ln a$ & $\dfrac{(1+x)^{p}-1}{x}\to p$ & $\dfrac{\tanh x}{x}\to1$ \\
\bottomrule
\end{tabular}
\end{table}

\section{Asíntotas}

\begin{definicion}[Asíntotas]\label{def:asintotas}
Sea $f$ una función.
\begin{enumerate}
  \item La recta $x=a$ es una \emph{asíntota vertical} si alguno de los límites laterales de $f$ en $a$ es $+\infty$ o $-\infty$.
  \item La recta $y=L$ es una \emph{asíntota horizontal} en $+\infty$ si $\lim_{x\to+\infty}f(x)=L$ (análogamente en $-\infty$).
  \item La recta $y=mx+b$ con $m\ne0$ es una \emph{asíntota oblicua} en $+\infty$ si $\lim_{x\to+\infty}\bigl(f(x)-mx-b\bigr)=0$.
\end{enumerate}
\end{definicion}

\begin{proposicion}[Cálculo de la asíntota oblicua]\label{prop:oblicua}
$y=mx+b$ es asíntota oblicua de $f$ en $+\infty$ si y solo si $m=\lim_{x\to+\infty}f(x)/x$ (finito y no nulo) y $b=\lim_{x\to+\infty}\bigl(f(x)-mx\bigr)$ (finito).
\end{proposicion}

\addtocounter{definicion}{1}%

\addtocounter{definicion}{1}%

\addtocounter{definicion}{1}%

